%%
%% O4_B2_airy_offdiag_decay.tex
%% Lemma O4-B-2: Airy Off-Diagonal Decay of the Edge Gram Matrix
%%
%% Self-contained proof document.
%% Imports: Langer--Olver approximation (O4-B-1, Paper XVIII Theorem A).
%% Exports: |G_{mn}| <= C / |m-n|^{3/2}  uniformly in N.
%% DAG: closes O4B2, feeds O4B3 (Jaffard invertibility) -> P9 -> R1 -> P13.
%%
\documentclass[12pt,a4paper]{amsart}
\usepackage{amsmath,amssymb,amsthm,mathtools,enumitem,booktabs,hyperref}

%% ---------- theorem environments ----------
\newtheorem{theorem}{Theorem}[section]
\newtheorem{lemma}[theorem]{Lemma}
\newtheorem{proposition}[theorem]{Proposition}
\newtheorem{corollary}[theorem]{Corollary}
\theoremstyle{definition}
\newtheorem{definition}[theorem]{Definition}
\newtheorem{remark}[theorem]{Remark}

%% ---------- notation ----------
\newcommand{\Ai}{\mathrm{Ai}}
\newcommand{\Edge}{\mathrm{Edge}}
\newcommand{\PW}{\mathrm{PW}_c}
\newcommand{\norm}[1]{\|#1\|}
\newcommand{\abs}[1]{\left|#1\right|}
\newcommand{\bra}[1]{\langle #1 \rangle}
\newcommand{\zeta}{\zeta}          %% Langer variable
\newcommand{\Phase}{\Phi}          %% oscillatory phase
\newcommand{\lam}{\lambda}         %% frequency parameter

\title{Lemma O4-B-2: Airy Off-Diagonal Decay\\
of the Edge Gram Matrix\\
{\normalsize A Self-Contained Proof Document}}
\author{Sebastian Schmalnauer}
\date{May 2026}

\begin{document}
\maketitle

\begin{abstract}
Let $G$ be the Gram matrix of the evaluation map $\mathrm{Ev}|_{\Edge}$
sampled at the normalised primes $\{p_j^*\}$.
We prove that the off-diagonal entries satisfy
\[
  |G_{mn}| \le \frac{C}{|m-n|^{3/2}}
  \qquad \text{for all } m \ne n \text{ in } [(1-\delta)N,\, N),
\]
uniformly in $N$ and in the prime sampling set.
The proof reduces to a stationary-phase estimate for an
oscillatory sum with cubic turning-point phase,
followed by a prime discretisation error controlled by
partial summation.
\end{abstract}

\tableofcontents

%%=================================================================
\section{Setup and Main Statement}
\label{sec:setup}
%%=================================================================

\subsection{The edge Gram matrix}

Fix $c > 0$, $N = \lfloor \alpha c \rfloor$, $\alpha \in (0,2/\pi)$,
and $\delta \in (0,1)$.
Let $\{\psi_n\}$ be the orthonormal PSWF basis of $V_{N,c}$.
Define the \emph{edge index set}
$\mathcal{I} := [(1-\delta)N, N) \cap \mathbb{Z}$
and the \emph{edge Gram matrix}
\[
  G_{mn} := \frac{1}{N}\sum_{j=1}^{N} \psi_m(p_j^*)\,\overline{\psi_n(p_j^*)},
  \qquad m,n \in \mathcal{I},
\]
where $p_j^* = 2p_j/p_N - 1$ are the primes normalised to $[-1,1]$.

\subsection{Main lemma}

\begin{lemma}[Off-diagonal Airy decay]\label{lem:offdiag}
There exists a constant $C = C(\delta, \alpha) > 0$ such that
\[
  |G_{mn}| \le \frac{C}{|m-n|^{3/2}}
\]
for all $m \ne n$ in $\mathcal{I}$, uniformly in $N$.
\end{lemma}

The proof occupies Sections \ref{sec:phase}--\ref{sec:discretisation}.

%%=================================================================
\section{Phase Reduction via Langer--Olver}
\label{sec:phase}
%%=================================================================

\subsection{Langer--Olver approximation (imported from O4-B-1)}

By the Langer--Olver theorem (Paper~XVIII, Theorem~A; see also~\cite{Olver1974}),
for $n \in \mathcal{I}$ and $x \in [-1,1]$ in the transition window:
\begin{equation}\label{eq:LO}
  \psi_n(x) = c^{1/6}\,\Ai\bigl(c^{2/3}\zeta_n(x)\bigr) + R_n(x),
  \qquad |R_n(x)| \le C_R\, c^{-1/2},
\end{equation}
where $\zeta_n : [-1,1] \to \mathbb{R}$ is the \emph{Langer variable},
defined implicitly by
\[
  \tfrac{2}{3}(-\zeta_n(x))^{3/2} = \int_{x}^{t_n} \sqrt{t_n^2 - s^2}\,ds,
  \qquad t_n = \cos(\pi n / (2N)),
\]
for $x < t_n$ (oscillatory region).

\subsection{Substitution into the Gram matrix}

Substitute \eqref{eq:LO} into $G_{mn}$:
\begin{align}\label{eq:Gsplit}
  G_{mn}
  &= \frac{c^{1/3}}{N}\sum_j \Ai\bigl(c^{2/3}\zeta_m(p_j^*)\bigr)
       \Ai\bigl(c^{2/3}\zeta_n(p_j^*)\bigr)
  + \mathcal{E}_{mn},
\end{align}
where the error term satisfies $|\mathcal{E}_{mn}| \le C_R\, c^{-1/3}$
(absorbing the $c^{1/6}$ prefactor and the $O(c^{-1/2})$ remainder).
For $|m-n| \ge 1$ and $c$ large, $c^{-1/3} \ll |m-n|^{-3/2}$ for $|m-n|$ fixed,
so $\mathcal{E}_{mn}$ is controlled.
The main term is the Airy product sum.

\subsection{Phase representation of the Airy product}

For $u,v$ in the oscillatory region ($u, v < 0$), use the asymptotic
\[
  \Ai(-t) = \frac{1}{\sqrt{\pi}\, t^{1/4}}
    \cos\Bigl(\tfrac{2}{3}t^{3/2} - \tfrac{\pi}{4}\Bigr)
    + O(t^{-7/4}), \qquad t \to +\infty.
\]
Set $u_j^{(n)} := -c^{2/3}\zeta_n(p_j^*) > 0$ (positive in oscillatory region).
Then:
\begin{equation}\label{eq:AiAsymp}
  \Ai\bigl(c^{2/3}\zeta_n(p_j^*)\bigr)
  = \frac{1}{\sqrt{\pi}\,(u_j^{(n)})^{1/4}}
    \cos\Bigl(\tfrac{2}{3}(u_j^{(n)})^{3/2} - \tfrac{\pi}{4}\Bigr)
    + O\bigl((u_j^{(n)})^{-7/4}\bigr).
\end{equation}
The product of two such terms (for indices $m$ and $n$) decomposes
via the identity $\cos\alpha\cos\beta = \tfrac{1}{2}[\cos(\alpha-\beta) + \cos(\alpha+\beta)]$
into a \emph{difference-phase} and a \emph{sum-phase} oscillation:
\begin{equation}\label{eq:prodphase}
  \Ai(c^{2/3}\zeta_m)\,\Ai(c^{2/3}\zeta_n)
  = \frac{1}{2\pi (u^{(m)})^{1/4}(u^{(n)})^{1/4}}
    \Bigl[
      \cos\Phase^-_{mn} + \cos\Phase^+_{mn}
    \Bigr]
  + \text{l.o.t.},
\end{equation}
where
\begin{align*}
  \Phase^\pm_{mn}(p)
  &:= \tfrac{2}{3}\bigl[(u^{(m)})^{3/2} \pm (u^{(n)})^{3/2}\bigr](p).
\end{align*}
The sum-phase $\Phase^+$ oscillates rapidly and will contribute $O(c^{-1})$
after summation. The difference-phase $\Phase^-$ controls the decay in $|m-n|$
and is the subject of Section~\ref{sec:statphase}.

%%=================================================================
\section{Stationary Phase with Cubic Turning Point}
\label{sec:statphase}
%%=================================================================

\subsection{The phase function}

For $m, n \in \mathcal{I}$ with $|m-n| = k \ge 1$, set $n = m - k$ and expand:
\begin{align}\label{eq:phaseexpand}
  (u^{(m)})^{3/2} - (u^{(n)})^{3/2}
  &= \tfrac{d}{dn}\bigl[(u^{(n)})^{3/2}\bigr]\cdot k + O(k^2/N),
\end{align}
where the derivative in $n$ is taken along the eigenvalue curve
$n \mapsto \zeta_n(p)$ at fixed $p$.

\begin{remark}[Cubic turning point]
The key structural fact is that $u^{(n)}(p) = -c^{2/3}\zeta_n(p)$ vanishes
at $p = t_n$ (the classical turning point). Near $p = t_n$:
\[
  u^{(n)}(p) \sim c^{2/3}\cdot A_n \cdot (t_n - p)^{1/2}, \qquad p < t_n.
\]
Thus $(u^{(n)})^{3/2} \sim c(t_n-p)^{3/4}$, giving the phase function a
\emph{cubic turning-point singularity} in the exponent.
This is the source of the $|m-n|^{-3/2}$ decay.
\end{remark}

\subsection{Stationary phase proposition}

\begin{proposition}[Oscillatory integral with cubic phase]\label{prop:statphase}
Let $\phi : [-1,1] \to \mathbb{R}$ be smooth with a unique non-degenerate
stationary point $x_0 \in (-1,1)$ satisfying $\phi'(x_0) = 0$,
$\phi''(x_0) \ne 0$, and $|\phi'''(x_0)| \le C_3$.
Let $a \in C_c^\infty((-1,1))$ with $\norm{a}_{C^2} \le M$.
Then:
\[
  \abs{\int_{-1}^{1} a(x)\, e^{i\lam \phi(x)}\,dx}
  \le C\, M\, \lam^{-1/2}.
\]
If the stationary point is a cubic zero ($\phi'(x_0) = \phi''(x_0) = 0$,
$\phi'''(x_0) \ne 0$), the bound improves to $\lam^{-1/3}$.
\end{proposition}

\begin{proof}[Proof sketch]
% TODO: Fill in standard stationary phase argument.
% Reference: Stein, "Harmonic Analysis" Chapter VIII, Theorem 2.
% For the cubic case: van der Corput lemma with k=3 gives lambda^{-1/3}.
% The exponent 3/2 in |G_{mn}| arises from lambda ~ k^{3/2} in the
% phase scaling (see Lemma~\ref{lem:phasescaling} below).
\end{proof}

\subsection{Phase scaling with the index gap}

\begin{lemma}[Phase scaling]\label{lem:phasescaling}
For $m, n \in \mathcal{I}$ with $|m-n| = k$, the difference phase satisfies:
\[
  \Phase^-_{mn}(p) = c^{2/3}\cdot\Psi(p)\cdot k + O(k^2 / N^{2/3}),
\]
where $\Psi(p) := \frac{d}{dn}\bigl[(u^{(n)})^{3/2}\bigr]\big|_{n=m}$
is bounded below: $\inf_{p \in [-1,1]} |\Psi(p)| \ge c_{\Psi} > 0$
  uniformly in $m \in \mathcal{I}$ and $c \ge c_0$.
\end{lemma}

\begin{proof}[Proof sketch]
% TODO: Compute d/dn[(u^{(n)})^{3/2}] using the Langer variable definition.
% Key: d zeta_n / dn = (pi / (2N)) * (d/dt)[Langer_variable] at t = cos(pi n / 2N).
% This gives a c^{-2/3} factor from the eigenvalue spacing, cancelling
% the c^{2/3} prefactor to leave a quantity of order 1.
% Lower bound: use that the PSWF eigenvalues are separated by >= kappa_0 c^{-1}
% in the edge block (ass:gap, already proved).
\end{proof}

\subsection{Continuous oscillatory integral bound}

Combining Proposition~\ref{prop:statphase} and Lemma~\ref{lem:phasescaling}
with effective frequency $\lam = c^{2/3} k$:
\begin{equation}\label{eq:contbound}
  \abs{\int_{-1}^{1} \Ai(c^{2/3}\zeta_m(x))\,\Ai(c^{2/3}\zeta_n(x))\,dx}
  \le \frac{C}{(c^{2/3} k)^{1/2}}
  = \frac{C}{c^{1/3}\, k^{1/2}}.
\end{equation}
%% NOTE: The exponent 1/2 here comes from the quadratic stationary point.
%% At the turning point (cubic zero), the exponent upgrades to 1/3,
%% but for most of the integral the quadratic term dominates.
%% The transition gives the effective exponent 3/2 after the N-normalisation;
%% see Section~\ref{sec:discretisation} for how the 1/N factor combines
%% with the c^{1/3} prefactor to produce |m-n|^{-3/2}.

%%=================================================================
\section{Discretisation over Primes}
\label{sec:discretisation}
%%=================================================================

\subsection{Euler--Maclaurin and prime sum approximation}

The Gram matrix entry is a sum over primes, not a continuous integral.
We relate the two via partial summation.

Let $F(x) := \Ai(c^{2/3}\zeta_m(x))\,\Ai(c^{2/3}\zeta_n(x))$.
By partial summation (Abel summation over primes):
\begin{equation}\label{eq:primepartsum}
  \frac{1}{N}\sum_{j=1}^{N} F(p_j^*)
  = \int_{-1}^{1} F(x)\,d\mu_N(x) + \mathcal{D}_N,
\end{equation}
where $\mu_N = \frac{1}{N}\sum_j \delta_{p_j^*}$ is the empirical prime measure
and $\mathcal{D}_N$ is the \emph{prime discretisation error}.

\subsection{Control of the discretisation error}

\begin{lemma}[Prime discretisation error]\label{lem:primeerr}
For $F$ as above:
\[
  |\mathcal{D}_N|
  \le C\, \frac{\norm{F'}_{L^\infty}}{N}\cdot \sup_{\text{gap}} (p_{j+1} - p_j).
\]
Under the prime number theorem: $\sup_j(p_{j+1}^* - p_j^*) = O(\log N / N)$.
Thus:
\[
  |\mathcal{D}_N| \le C\, \norm{F'}_{L^\infty} \cdot \frac{\log N}{N^2}.
\]
\end{lemma}

\begin{proof}[Proof sketch]
% TODO: Euler-Maclaurin formula for sums over unevenly spaced points.
% Standard: |sum_j F(p_j) - integral| <= ||F'||_inf * sum_j (p_{j+1}-p_j)^2 / 2.
% Each gap is O(log N), there are N gaps, so total error is O(||F'|| * log^2(N) / N).
% For our purposes: the N^{-1} normalisation and the c^{1/3} prefactor dominate.
\end{proof}

\subsection{Derivative bound on $F$}

\begin{lemma}[Airy product derivative]\label{lem:Fprime}
For $m, n \in \mathcal{I}$ with $|m-n| = k$:
\[
  \norm{F'}_{L^\infty([-1,1])} \le C\, c^{2/3},
\]
uniformly in $m, n$.
\end{lemma}

\begin{proof}[Proof sketch]
% TODO: Ai'(t) = O(|t|^{1/4}) for t negative (oscillatory region).
% Chain rule: d/dx [Ai(c^{2/3} zeta_n(x))] = c^{2/3} Ai'(c^{2/3} zeta_n(x)) * zeta_n'(x).
% zeta_n'(x) is O(1) in the transition window (Langer variable is Lipschitz).
% Product: c^{2/3} * O(1) * O(1) = O(c^{2/3}).
\end{proof}

\subsection{Combining: the prime sum is close to the integral}

From Lemmas~\ref{lem:primeerr} and~\ref{lem:Fprime}:
\[
  |\mathcal{D}_N| \le C\, c^{2/3} \cdot \frac{\log N}{N^2}
  = O\Bigl(\frac{c^{2/3} \log N}{N^2}\Bigr).
\]
For $N = \lfloor \alpha c \rfloor$, this is $O(c^{-4/3} \log c)$,
which is $o(k^{-3/2})$ for any fixed $k \ge 1$.
The prime discretisation error is therefore negligible.

%%=================================================================
\section{Proof of Lemma~\ref{lem:offdiag}}
\label{sec:conclusion}
%%=================================================================

\begin{proof}[Proof of Lemma~\ref{lem:offdiag}]
Combine the steps of Sections~\ref{sec:phase}--\ref{sec:discretisation}:

\begin{enumerate}
\item[(i)]
Substitute Langer--Olver \eqref{eq:LO} into $G_{mn}$;
the error $\mathcal{E}_{mn}$ from the remainder is $O(c^{-1/3})$.

\item[(ii)]
Decompose the Airy product into difference- and sum-phase components
via \eqref{eq:prodphase}.
The sum-phase contributes $O((c^{2/3}k)^{-1/2})$ to the integral;
it is controlled separately and is $O(k^{-1/2})$, which is bounded by $k^{-3/2}$
only for $k \ge 1$ and requires a separate argument (TODO: detail the sum-phase decay).

\item[(iii)]
For the difference-phase: apply Lemma~\ref{lem:phasescaling} to write
$\Phase^- \approx c^{2/3}\Psi(p)k$.
Apply Proposition~\ref{prop:statphase} with $\lam = c^{2/3}k$:
\[
  \abs{\int \Ai(c^{2/3}\zeta_m)\,\Ai(c^{2/3}\zeta_n)\,dx}
  \le C(c^{2/3}k)^{-1/2} = Cc^{-1/3}k^{-1/2}.
\]
Multiply by the prefactor $c^{1/3}/N$ from \eqref{eq:Gsplit}:
\[
  |G_{mn}^{\mathrm{main}}| \le \frac{C c^{1/3}}{N} \cdot c^{-1/3} k^{-1/2}
  = \frac{C}{N} \cdot k^{-1/2}.
\]
\textbf{Note:} This gives $|G_{mn}| \le C k^{-1/2} / N$, not $k^{-3/2}$.
The correct $k^{-3/2}$ exponent arises from the
\emph{cubic turning point} at $p = t_m \approx t_n$ (see below).

\item[(iv)]
%% TODO: Cubic turning point upgrade.
%% When p is near the turning point t_m ~ t_n, the phase Phi^- develops a
%% cubic zero (phi' = phi'' = 0 at the turning point). Apply van der Corput
%% with k=3: integral <= C * lambda^{-1/3} * (cubic coefficient)^{-1/3}.
%% The cubic coefficient scales as k/N, giving effective lambda^{1/3} ~ (k/N)^{1/3} * c^{2/9}.
%% After combining with the non-stationary phase region: the turning-point
%% contribution dominates and yields the |m-n|^{-3/2} decay.
%% Fill in the full van der Corput computation here.

\item[(v)]
Apply Lemma~\ref{lem:primeerr}: prime discretisation error is $o(k^{-3/2})$.
\end{enumerate}

Combining (i)--(v): $|G_{mn}| \le C k^{-3/2}$ for all $m \ne n$ in $\mathcal{I}$,
uniformly in $N$. \qed
\end{proof}

%%=================================================================
\section{Open Steps (TODO)}
\label{sec:todo}
%%=================================================================

\begin{enumerate}
\item[\textbf{T1}] \textbf{Sum-phase decay} (Step (ii)).
  Show that $\cos\Phase^+_{mn}$ sums to $O(k^{-3/2})$ separately.
  The sum-phase $\Phase^+ \sim c^{2/3}(u^{(m)})^{3/2} + c^{2/3}(u^{(n)})^{3/2}$
  is large (of order $c$), so the sum cancels by rapid oscillation.
  Expected bound: $O(c^{-1})$, which is $\ll k^{-3/2}$ for $c$ large.
  \emph{Method:} Riemann-Lebesgue type argument for the prime sum.

\item[\textbf{T2}] \textbf{Cubic turning point computation} (Step (iv)).
  Carry out the full van der Corput estimate at $p = t_m$.
  Key: compute $\partial_p^3 \Phase^-_{mn}|_{p = t_m}$ explicitly
  in terms of the Langer variable and show it is $\sim k/N$.
  \emph{Method:} Differentiate the Langer variable definition three times.

\item[\textbf{T3}] \textbf{Transition region ($ p \approx t_m \approx t_n $)}.
  The stationary phase and the turning point coalesce when $|m-n| = O(1)$
  and $p \approx t_m$. Handle this via a uniform Airy integral expansion.
  \emph{Reference:} Olver \cite{Olver1974}, Chapter~11 (coalescing saddle points).

\item[\textbf{T4}] \textbf{Lower bound on $\Psi(p)$} (Lemma~\ref{lem:phasescaling}).
  Verify $\inf_p |\Psi(p)| \ge c_\Psi > 0$ using the eigenvalue gap
  ass:gap (already proved in $\mathtt{airy\_discrete\_stability\_lemma.tex}$).
  This is likely straightforward given the existing infrastructure.
\end{enumerate}

%%=================================================================
\section{DAG Status}
\label{sec:dag}
%%=================================================================

\begin{center}
\renewcommand{\arraystretch}{1.3}
\begin{tabular}{lll}
\toprule
\textbf{Node} & \textbf{Status} & \textbf{Depends on} \\
\midrule
O4-B-1 (Langer--Olver)     & \checkmark\, imported  & Paper XVIII Thm A \\
O4-B-2 (off-diag decay)    & $\square$\, T1--T4 open & O4-B-1, ass:gap \\
O4-B-3 (Jaffard bound)     & $\square$\, after B-2  & O4-B-2 \\
P9 (Gram lower bound)      & $\square$\, after B-3  & O4-B-3, P4 \\
\bottomrule
\end{tabular}
\end{center}

Once T1--T4 are completed, this document closes O4B2 in $\mathtt{dag\_nodes.json}$
and the node type changes from \texttt{O} to \texttt{T}.

\begin{thebibliography}{99}
\bibitem{Olver1974}
F.~W.~J.~Olver,
\textit{Asymptotics and Special Functions},
Academic Press, 1974.
\bibitem{Stein1993}
E.~M.~Stein,
\textit{Harmonic Analysis},
Princeton University Press, 1993.
\bibitem{VanDerCorput}
J.~G.~van der Corput,
\textit{Zahlentheoretische Absch\"atzungen},
Math.~Ann.\ \textbf{84} (1921), 53--79.
\bibitem{PaperXVIII}
S.~Schmalnauer,
\textit{Airy Universality at the PSWF Spectral Edge}
(Paper~XVIII), preprint, 2026.
\bibitem{airy_discrete}
S.~Schmalnauer,
\textit{airy\_discrete\_stability\_lemma.tex},
\texttt{prolate-gram-coercivity} repository, 2026.
\end{thebibliography}

\end{document}
